\chapter{Despegue: carrera y fase aérea}\label{ch:takeoff}
\objetivos{Integrar la aceleración sobre pista; distinguir velocidad aérea y terrestre; añadir rotación y obstáculo con hipótesis explícitas; separar cálculo académico y actuación certificada.}

\section{Fases y magnitudes}
La distancia de despegue de un modelo puede separarse en carrera, rotación, transición y ascenso hasta una altura de referencia. La definición de cada fase debe mantenerse al comparar resultados. Una distancia hasta separación de ruedas no equivale a distancia hasta obstáculo; una distancia calculada con empuje constante no equivale a una distancia certificada con procedimientos, fallos y márgenes reglamentarios.

Se estudia una pista horizontal, viento alineado, avión rígido y una ley de empuje finita. Sea $V_g\geq0$ la velocidad terrestre y $U>0$ viento de cara, de modo que $V_a=V_g+U$ mientras el flujo conserve su sentido. La distancia satisface $\dot x=V_g$; la sustentación y resistencia se calculan con $V_a$.

\section{Balance en carrera}
Sin aceleración vertical y antes de perder contacto,
\begin{equation}
N=W-L,\qquad
m\dot V_g=T-D-\mu_r(W-L),\qquad N\geq0.
\label{eq:takeoffroll}
\end{equation}
La resistencia de rodadura se aproxima como $\mu_rN$. Una pendiente de pista añade una componente de peso y modifica la reacción normal. La expresión debe cambiar si se incluyen fuerzas verticales de propulsión o amortiguación del tren.

Para una actitud de carrera fija se pueden adoptar $C_{L,g}$ y $C_{D,g}=C_{D0}+kC_{L,g}^2$. No se utiliza $L=W$ durante toda la carrera: al principio las ruedas soportan casi todo el peso y la sustentación crece con velocidad aérea.

\section{Integración en velocidad}
Definiendo $a(V_g)=[T(V_a)-D(V_a)-\mu_r(W-L(V_a))]/m$,
\begin{equation}
s_g=\int_0^{V_{g,\rm fin}}\frac{V_g}{a(V_g)}\,\dd V_g,
\qquad t_g=\int_0^{V_{g,\rm fin}}\frac{1}{a(V_g)}\,\dd V_g.
\label{eq:rollintegral}
\end{equation}
Se exige $a>0$ en el intervalo para una carrera monótona. Si aparece un punto donde $a=0$ antes de la velocidad objetivo, ésta no se alcanza en el modelo. Una cuadratura que atraviese esa singularidad y devuelva un valor no describe una carrera posible.

\subsection{Solución cerrada sin viento y empuje constante}
Si $T$, $C_{L,g}$ y $\mu_r$ son constantes,
\begin{equation}
a(V)=a_0+a_2V^2,
\quad a_0=T/m-\mu_rg_0,
\quad a_2=\frac{\rho S}{2m}(\mu_rC_{L,g}-C_{D,g}).
\end{equation}
Para $a_2\neq0$,
\begin{equation}
s_g=\frac{1}{2a_2}\ln\!\left(\frac{a_0+a_2V_{\rm fin}^2}{a_0}\right).
\label{eq:rolllog}
\end{equation}
El argumento debe ser positivo y la aceleración mantenerse positiva. En el límite $a_2\to0$ se recupera $V_{\rm fin}^2/(2a_0)$. Esta comprobación es útil para evitar inestabilidad numérica por dividir un logaritmo pequeño entre un número pequeño; una implementación puede usar una función de logaritmo estable o el límite.

El tiempo se obtiene con arcotangente si $a_2>0$ y con arcotangente hiperbólica si $a_2<0$, dentro del dominio de aceleración positiva. La cuadratura directa de $1/a$ evita confundir ambos casos y permite mapas de motor más realistas.

\section{Velocidad de separación y rotación}
Una condición académica posible es $V_{a,\rm LOF}=k_VV_{\rm stall,TO}$ con $k_V>1$ dado por el problema. No se presenta como regla operativa universal. Con viento de cara, la velocidad terrestre objetivo es $V_{g,\rm LOF}=V_{a,\rm LOF}-U$, si es positiva.

La rotación cambia incidencia y sustentación y requiere momento de elevador. Un modelo simplificado puede añadir una duración de rotación $t_R$ y distancia aproximada $V_{g,R}t_R$; debe aclarar que no se han integrado las ecuaciones de cabeceo. Un modelo de seis grados de libertad puede incluir esa dinámica, el tren y sus reacciones, pero necesita datos y validación apropiados.

\section{Transición y obstáculo}
Si se idealiza una transición circular a velocidad constante y radio $R_c$, desde trayectoria horizontal hasta $\gamma_c>0$,
\begin{equation}
x_{\rm tr}=R_c\sin\gamma_c,\qquad
h_{\rm tr}=R_c(1-\cos\gamma_c).
\end{equation}
En el inicio, $R_c\simeq V^2/[g_0(n-1)]$ para un factor de carga $n>1$. Conservar ese radio y carga durante toda la transición es una aproximación, porque la componente normal del peso cambia con el ángulo. Para un obstáculo $H>h_{\rm tr}$ y ascenso recto posterior,
\begin{equation}
x_{\rm asc}=\frac{H-h_{\rm tr}}{\tan\gamma_c}.
\end{equation}
La distancia total combina carrera, rotación, transición y tramo recto según la definición elegida. Si el obstáculo se alcanza dentro de la transición, no se añade un tramo recto de longitud negativa: se resuelve el ángulo de la circunferencia a esa altura.

\section{Masa, temperatura y viento}
Mayor masa eleva la velocidad aérea requerida y reduce aceleración para la misma fuerza. Menor densidad eleva TAS de separación y puede reducir empuje o potencia disponibles. Temperatura alta modifica densidad y velocidad del sonido; no se reduce a sumar metros de altitud sin declarar el modelo.

El viento de cara reduce velocidad terrestre necesaria para obtener la misma velocidad aérea y suele acortar carrera. También altera fuerzas durante la carrera, porque ya existe presión dinámica al inicio. El viento de cola puede producir un intervalo de flujo inverso a velocidad terrestre pequeña: un modelo que sólo admite $V_a\geq0$ debe rechazarlo o incorporar la aerodinámica correspondiente.

\ilustra{despegue-masa}{Distancia de rodadura en función de masa con una ley de empuje constante y velocidad final ligada a pérdida. No incluye rotación, obstáculo ni márgenes de certificación.}

\caso{Ejemplo original y dos métodos}{Con $m=1.500\SIu{kg}$, $S=18\SIu{m^2}$, $\rho=1.225\SIu{kg\,m^{-3}}$, polar $0.035+0.055C_L^2$, $C_{L,g}=0.65$, $\mu_r=0.03$, $T=5.000\SIu N$ y objetivo terrestre $32\SIu{m\,s^{-1}}$ sin viento, se integra la ecuación~\eqref{eq:rollintegral}. Se comprueba que $N>0$ hasta el final y que $a>0$. La cuadratura se contrasta con una marcha temporal con pasos 0.02, 0.01 y 0.005 s. Los resultados pasan un error relativo de distancia de $2\times10^{-4}$ respecto a la cuadratura. La velocidad final se ha prescrito: no se afirma que sea la velocidad certificada de separación.}

\section{Límites de una actuación académica}
La distancia calculada depende de rozamiento, ley de empuje, técnica de rotación, configuración, superficie y ambiente. Las actuaciones certificadas añaden definiciones, condiciones y márgenes específicos, y se usan mediante los documentos aprobados de la aeronave. Este capítulo enseña a construir y auditar un modelo, sin transformar su resultado en una autorización de operación.

\section{Taller resuelto: Viento de cara: fuerza y distancia}

\noindent\textbf{Nivel:} Intermedio.

\apartado{1. Enunciado}{Para masa 1680 kg, S=20.2 m², rho=1.20 y empuje 5400 N, calcule carrera hasta TAS 34 m/s en calma y con viento de cara 5 m/s. Use rodadura 0.03, CL de carrera 0.63 y polar 0.034+0.051CL².}

\apartado{2. Datos}{La velocidad terrestre objetivo cambia de 34 a 29 m/s.}

\apartado{3. Incógnitas}{Las magnitudes y comprobaciones solicitadas en el caso.}

\apartado{4. Sistema de referencia}{Pista horizontal; fuerzas con velocidad aérea y distancia con velocidad terrestre.}

\apartado{5. Hipótesis físicas}{Empuje finito constante; sin rotación ni obstáculo; mismo estado atmosférico y configuración.}

\apartado{6. Modelo matemático}{Viento de cara: fuerza y distancia}

\apartado{7. Ecuaciones simbólicas}{Se parte de los balances y definiciones siguientes.}

\begin{equation}V_a=V_g+U\end{equation}

\begin{equation}a=[T-D(V_a)-\mu(W-L(V_a))]/m\end{equation}

\begin{equation}s=\int_0^{V_{g,f}}V_g/a(V_g)\,\dd V_g\end{equation}

\apartado{8. Desarrollo algebraico}{Las operaciones se realizan en este orden, conservando las condiciones del modelo.}

\begin{enumerate}

\item Se prepara un caso en calma y otro con su objetivo terrestre distinto.

\item En el caso con viento las fuerzas se evalúan con velocidad aérea desde Vg=0.

\item Se comprueban aceleración positiva y contacto y se integran ambas carreras.

\item Se compara reducción de distancia y se distingue de una simple corrección de velocidad final.

\end{enumerate}

\apartado{9. Sustitución numérica}{Se usan las magnitudes del apartado de datos en SI; los ángulos se convierten a radianes antes de evaluar funciones trigonométricas. Las densidades indicadas como ISA se obtienen de la capa correspondiente.}

\noindent Carrera en calma:\quad $208.6565284657509\simeq 208.6565\,\mathrm m$.\par

\noindent Carrera con viento:\quad $152.6346719423097\simeq 152.6347\,\mathrm m$.\par

\noindent Reducción relativa:\quad $0.2684883954284547\simeq 0.2684884\,$.\par

\apartado{10. Resultados con unidades}{Los resultados pertenecen al modelo y al intervalo de datos declarados.}

\begin{center}\begin{tabular}{lr}\toprule Magnitud&Resultado\\\midrule

Carrera en calma & $208.6565\,\mathrm m$\\

Carrera con viento & $152.6347\,\mathrm m$\\

Reducción relativa & $0.2684884\,$\\

\bottomrule\end{tabular}\end{center}

\apartado{11. Interpretación}{El viento reduce la carrera por varios términos del mismo balance; no es sólo una resta de avance al final.}

\apartado{12. Comprobación independiente y alcance}{Marcha temporal con eventos de velocidad aérea y revisión de fuerzas iniciales en ambos casos.}

\apartado{13. Limitaciones}{No se estudia viento de cola con flujo invertido ni ráfagas.}

\par\smallskip\noindent{\footnotesize\textit{Datos y solución originales. Las sustituciones registradas se reproducen con SymPy y Studium; eso no establece validación experimental ni revisión experta.}}


\repaso{¿Por qué se integra $V_g/a$ y no $V_a/a$ para obtener pista?}{Porque la pista se mide respecto al suelo y $\dot x=V_g$. Las fuerzas se calculan con velocidad aérea. Sustituir una por otra en presencia de viento mezcla cinemática y aerodinámica.}
\resumen{La rodadura exige balances con reacción normal variable y una ley de empuje finita. Las fases aéreas se añaden con geometría y dinámica declaradas. Viento, masa y densidad afectan tanto a la velocidad objetivo como a la aceleración.}
\fuentes{Guía ULE, despegue. Familias de ecuaciones identificadas en el formulario facilitado y derivadas aquí desde Newton. El ejemplo es académico y original; no se atribuye a un procedimiento FAA/EASA ni a una aeronave certificada.}
